
Paper Session 1: Digital Futures within Limits

	How digital will the future be? Analysis of prospective scenarios
Aurélie Bugeau, Anne-Laure Ligozat
        https://arxiv.org/abs/2312.15948
        https://doi.org/10.48550/arXiv.2312.15948

	Computing, Degrowth and Complexity : Systemic Considerations for Digital De-escalation
Valentin Girard, Maud Rio, Romain Couillet

	Digital Sobriety: from Tips to Values
Nicolas Szilas


Paper Session 2: Limits in Education

	Towards designing symbiotic university-adjacent sustainability-focused computing cooperatives as edge spaces for navigating transition: an imaginary and initial feasibility analysis
Greg L. Nelson

	What do we Teach when Teaching More-than-human Perspectives in Computing and Technology Design Education? – An Emerging Pedagogical Framework
Elisabet M. Nilsson, Rikke Hagensby Jensen, Anne-Marie Hansen, Tilde Bekker, Daisy Yoo, Eva Eriksson

	Interaction design pedagogy for digital degrowth: lessons learned from student inquiry
Cathryn Ploehn, Carla Garcia Leija, Grace Park, Chloe Kim



Paper Session 3: Computing for the Energy Transition
	Windternet: Designing grid-liberated servers for regenerative energy communities
Eric Snodgrass, Helen Pritchard, Miranda Moss, Daniel Gustafsson, Jorge Luis Zapico

	Fake Dashboards Result in Fake Insights: The Challenges of Prototyping Energy Dashboards
Christina Bremer, Christian Remy, Adrian Friday

	Energy Experience Design
Brian Sutherland



Paper Session 4: Speculative Designs within Limits

	Limits at a Distance: Design Directions to Address Psychological Distance in Climate Policy Decisions
Eshta Bhardwaj, Han Qiao, Christoph Becker
        https://arxiv.org/abs/2507.18913
        https://doi.org/10.48550/arXiv.2507.18913

	Exploring post-neoliberal futures for managing commercial heating and cooling through speculative praxis
Oliver Bates, Christian Remy, Kieran Cutting, Adam Tyler, Adrian Friday
        https://arxiv.org/abs/2507.19072
        https://doi.org/10.48550/arXiv.2507.19072
        

